\documentclass[11pt,a4paper]{article}
\usepackage[T1]{fontenc}
\usepackage[latin1]{inputenc}
\usepackage{ngerman}
\usepackage{a4wide}
\usepackage[dvips]{graphicx}

\usepackage[
pdfauthor={},
pdftitle={System Requirements},
pdfcreator={pdftex},
]{hyperref}

\usepackage{sectsty}
\allsectionsfont{\sffamily}

\usepackage{fancyheadings} 
\pagestyle{fancy} 
\lhead{\textsf{\textbf{ACE} \\ \small{a collaborative editor}}}
\chead{}
\rhead{\includegraphics[height=0.875cm,width=3cm]{../../images/logo_BFH.eps}}
\lfoot{}
\cfoot{\textsf{\thepage}}
\rfoot{}
\setlength{\headrulewidth}{0.6pt}
\setlength{\footrulewidth}{0.6pt}
\setlength{\topmargin}{-50pt}
\addtolength{\headheight}{50pt}

\usepackage{colortbl}

\newcommand{\headercol}[2]{\multicolumn{1}{|>{\bfseries\columncolor[gray]{0.82}}p{#1}|}{\textsf{#2}}}
\newcommand{\ace}[0]{\emph{ACE }}



\begin{document}

\begin{titlepage}
\thispagestyle{empty}
  \includegraphics[height=1.5in]{../../images/pix.eps}

  \begin{center}

    {\fontsize{40}{45} \textbf{\textsf{ACE}}} \\
    \textsf{a collaborative editor} \\
        
    \vspace{36pt}
        
    {\huge{\textbf{\textsf{Report Evaluation Network}}}} \\

    \vspace{36pt}

	\textsf{Berne University of Applied Sciences} \\
    \textsf{School of Engineering and Information Technology} \\
    
  \end{center}

  \vfill
  
  \begin{tabular}{ll}
   \hline

   \\

   \multicolumn{1}{>{\bfseries}p{1.5in}}{\textsf{Date:}} &
   \multicolumn{1}{>{}p{4.3in}}{\textsf{14.06.2005}}          \\
   
   \\
   
   \multicolumn{1}{>{\bfseries}p{1.5in}}{\textsf{Version:}}     &   
   \multicolumn{1}{>{}p{4.3in}}{\textsf{0.6}}                 \\

   \\
   
   \multicolumn{1}{>{\bfseries}p{1.5in}}{\textsf{Projectteam:}}                 &
   \multicolumn{1}{>{}p{4.3in}}{\textsf{Mark Bigler (biglm2@hta-bi.bfh.ch)}}  \\
   \multicolumn{1}{>{\bfseries}p{1.5in}}{}                                      &
   \multicolumn{1}{>{}p{4.3in}}{\textsf{Simon R�ss (rasss@hta-bi.bfh.ch)}}    \\
   \multicolumn{1}{>{\bfseries}p{1.5in}}{}                                      &
   \multicolumn{1}{>{}p{4.3in}}{\textsf{Lukas Zbinden (zbinl@hta-bi.bfh.ch)}} \\   
   
   \\
   
   \multicolumn{1}{>{\bfseries}p{1.5in}}{\textsf{Receivers:}}                       &
   \multicolumn{1}{>{}p{4.3in}}{\textsf{Jean-Paul Dubois (doj@hta-bi.bfh.ch)}}       \\
   \multicolumn{1}{>{\bfseries}p{1.5in}}{}                                          &
   \multicolumn{1}{>{}p{4.3in}}{\textsf{Claude Fuhrer (frc@hta-bi.bfh.ch)}}       \\

   \\
   
   \multicolumn{1}{>{\bfseries}p{1.5in}}{\textsf{Location:}}               &   
   \multicolumn{1}{>{}p{4.3in}}{\textsf{Subversion Repository}} \\

   \\  
   
   \hline
  \end{tabular}

\end{titlepage}


\section{Introduction}

\section{Project Goals}

The goal of the project ACE is to create a collaborative text editor. In the following a more detailed
description of what the functionality the application must have are described.

\subsection{Graphical User Interface}

The following list shows the basic features of the text editor:

\begin{itemize}
 \item create new documents
 \item load and save documents
 \item open several documents at the same time
 \item cut/copy/paste
 \item about information
\end{itemize}

\subsection{Collaboration Functionality}

\begin{itemize}
 \item publish local documents
 \item automatic discovery of other running instances of ACE on a local area network
 \item automatic discovery of published documents of other running instances of ACE
 \item manual discovery of running instances outside of the local area network
 \item join published documents
 \item leave joined documents
 \item kick users
\end{itemize}

One of the core principles for ACE was that we wanted it to be as simple as possible to set up. Most users
will feel lost if they have to enter IP addresses and ports. So ACE has to support automatic discovery
of other users (other running instances of ACE) and their shared documents. These automatically discovered
users and documents must be displayed in the graphical user interface where the user has a set of actions
to interact with them.


\section{Optional Goals}

We might implement optional goals if we have additional time.

\begin{itemize}
 \item access right mechanism
 \item invite other users to a session 
\end{itemize}

\subsection{Access Right Mechanism}

There is no mandatory goal that requires any kind of access right mechanism meaning that any user can join a published
document and gain read/write access. Of course, this can be an undesirable situation, especially in an environement
where you do not trust all the other users on the local area network.

A access right mechanism should allow the publisher to choose default access rights for each of his published 
documents. Reasonable access right settings would be:

\begin{itemize}
 \item ask for join permission
 \item read-only access
 \item read-write access
\end{itemize}

The first access mode ensures that nobody can access or even view a shared document before the publisher explicitely gives
that user the permission to view the document, either with read-only access or full read-write access.


\section{Non-Goals}

Due to the limited available time, certain (useful) functionality cannot be implemented in the diploma project. The 
principle is that if something is not explicitely listed in the project goals, it is not a goal of the project. However,
we want to highlight in this section some points that are not goals.

\begin{itemize}
 \item security aspects: authentication, confidentiality, authorization
 \item undo/redo
\end{document}
